Partisan gerrymandering reduces electoral competitiveness by packing
opposition voters into supermajority districts and cracking them across
many hopeless districts.
Algorithmic redistricting cannot access partisan data and therefore cannot
deliberately create safe or competitive seats.
We test whether compact algorithmic plans nonetheless produce systematically
more competitive districts than enacted gerrymanders.

Using \textsc{bisect} 2020 congressional plans and enacted 2022 plans for
five study states (NC, WI, TX, FL, PA), we compute the fraction of competitive
districts (presidential vote margin $< 10$ percentage points) and the
efficiency gap using VEST 2020 presidential precinct returns
interpolated to district assignments.

\textsc{Bisect} plans produce a competitive district fraction of
$0.40 \pm 0.06$ across the five states, compared to $0.17 \pm 0.09$
for enacted gerrymanders.
The efficiency gap is $-2.9\%$ (slight Democratic advantage) for bisect plans
vs.\ $+5.8\%$ (Republican advantage) for enacted plans ---
a difference of 8.7 percentage points.
The competitive district increase is explained by two mechanisms:
(1) algorithmic plans cannot deliberately pack opposition voters,
so partisan composition within districts reflects the underlying geographic
distribution rather than mapmaker choices;
(2) compactness optimisation produces spatially contiguous districts that
mix partisan communities rather than reaching across geographic features
to concentrate them.

These results confirm that algorithmic redistricting is not merely
geometrically neutral --- it is electorally more competitive than the
partisan gerrymanders it would replace.
